% from ega1-10-fr.tex
\oldpage[I]{207}

un système inductif de $S_n$-morphismes $u_n:X_n\to Y_n$ tel que $u_m$
s'identifie à $(u_n)_{(S_m)}$ pour $m\leq n$.

Le fait qu'on a bien défini ainsi une équivalence résultera de la proposition
plus précise suivante~:

\begin{proposition}[10.12.3.1]
\label{I.10.12.3.1-fr}
Soit $(X_n)$ un système inductif de $S_n$-préschémas~; on suppose que les
morphismes structuraux $f_n:X_n\to S_n$ sont tels que les diagrammes
(\hyperref[I.10.12.2.1-fr]{10.12.2.1}) soient commutatifs et identifient
$X_m$ à $X_n\times_{S_n}S_m$ pour $m\leq n$. Alors le système inductif
$(X_n)$ vérifie les conditions b) et c) de
(\hyperref[I.10.6.3-fr]{10.6.3})~; soient $\mathfrak{X}$ sa limite inductive,
$f:\mathfrak{X}\to\mathfrak{S}$ le morphisme limite inductive du système
inductif $(f_n)$. Alors, si $X_0$ est localement noethérien,
$\mathfrak{X}$ est localement noethérien et $f$ est un morphisme adique.
\end{proposition}

Comme le faisceau d'idéaux de $\mathscr{O}_{S_n}$ qui définit le
sous-préschéma $S_m$ de $S_n$ est nilpotent, il en est de même, en vertu de
(\hyperref[I.4.4.5-fr]{4.4.5}) du faisceau d'idéaux de
$\mathscr{O}_{X_n}$ définissant le sous-préschéma $X_m$ de $X_n$, donc les
conditions de (\hyperref[I.10.6.3-fr]{10.6.3}) sont bien vérifiées. La
question est par suite locale sur $\mathfrak{X}$ et $\mathfrak{S}$ et on
peut supposer que $\mathfrak{S}=\operatorname{Spf}(A)$,
$\mathscr{J}=\mathfrak{J}^\Delta$, $A$ étant un anneau $\mathfrak{J}$-adique
noethérien, et $X_n=\operatorname{Spec}(B_n)$~; si
$A_n=A/\mathfrak{J}^{n+1}$, l'hypothèse entraîne que $B_0$ est noethérien et
que si on pose $\mathfrak{J}_n=\mathfrak{J}/\mathfrak{J}^{n+1}$,
$B_m=B_n/\mathfrak{J}_n^{m+1}B_n$. Le noyau de $B_n\to B_0$ est donc
$\mathfrak{K}_n=\mathfrak{J}_nB_n$ et le noyau de $B_n\to B_m$ est
$\mathfrak{K}_n^{m+1}$ pour $m\leq n$~; en outre, comme $A_1$ est
noethérien, $\mathfrak{J}_1$ est de type fini sur $A_1$, donc
$\overline{\mathfrak{K}}_1=\mathfrak{K}_1/\mathfrak{K}_1^2$ est de type fini
sur $B_1$, et \emph{a fortiori} sur
$B_0=B_1/\mathfrak{K}_1$~; le fait que $\mathfrak{X}$ soit noethérien
résulte alors de (\hyperref[I.10.6.4-fr]{10.6.4})~; si
$B=\varprojlim B_n$, on a $\mathfrak{X}=\operatorname{Spf}(B)$, et si
$\mathfrak{K}$ est le noyau de $B\to B_0$,
$B_n=B/\mathfrak{K}^{n+1}$. Si
$\rho_n:A/\mathfrak{J}^{n+1}\to B/\mathfrak{K}^{n+1}$ est l'homomorphisme
correspondant à $f_n$, on a donc
\[
  \mathfrak{K}/\mathfrak{K}^{n+1}
  =(B/\mathfrak{K}^{n+1})\rho_n(\mathfrak{J}/\mathfrak{J}^{n+1})
\]
comme l'homomorphisme $\rho:A\to B$ correspondant à $f$ est égal à
$\varprojlim\rho_n$, l'idéal $\mathfrak{J}B$ de $B$ est dense dans
$\mathfrak{K}$, et comme tout idéal de $B$ est fermé
(\hyperref[0.7.3.5-fr]{0, 7.3.5}), on a $\mathfrak{K}=\mathfrak{J}B$. Si
$\mathscr{K}=\mathfrak{K}^\Delta$, la relation
$f^*(\mathscr{J})\mathscr{O}_{\mathfrak{X}}=\mathscr{K}$ résulte alors de
(\hyperref[I.10.10.9-fr]{10.10.9}) et achève la démonstration.

\begin{env}[10.12.3.2]
\label{I.10.12.3.2-fr}
L'équivalence précédente fournit, pour deux $\mathfrak{S}$-préschémas
adiques $\mathfrak{X}$, $\mathfrak{Y}$, une \emph{bijection canonique}
\[
  \operatorname{Hom}_{\mathfrak{S}}(\mathfrak{X},\mathfrak{Y})
  \xrightarrow{\sim}
  \varprojlim_n\operatorname{Hom}_{S_n}(X_n,Y_n)
\]
la limite projective étant relative aux applications
$u_n\to(u_n)_{(S_m)}$ pour $m\leq n$.
\end{env}

\subsection{Morphismes de type fini.}
\label{subsection:I.10.13-fr}

\begin{proposition}[10.13.1]
\label{I.10.13.1-fr}
Soient $\mathfrak{Y}$ un préschéma formel localement noethérien,
$\mathscr{K}$ un Idéal de définition de $\mathfrak{Y}$,
$f:\mathfrak{X}\to\mathfrak{Y}$ un morphisme de préschémas formels. Les
conditions suivantes sont équivalentes~:
\begin{enumerate}
  \item[a)] $\mathfrak{X}$ est localement noethérien, $f$ est un morphisme
    adique (\hyperref[I.10.12.1-fr]{10.12.1}) et si l'on pose
    $\mathscr{J}=f^*(\mathscr{K})\mathscr{O}_{\mathfrak{X}}$, le morphisme
    $f_0:(\mathfrak{X},\mathscr{O}_{\mathfrak{X}}/\mathscr{J})
    \to(\mathfrak{Y},\mathscr{O}_{\mathfrak{Y}}/\mathscr{K})$ déduit de
    $f$ est de type fini.
  \item[b)] $\mathfrak{X}$ est localement noethérien, et est limite
    inductive d'un $(Y_n)$-système inductif adique $(X_n)$ tel que le
    morphisme $X_0\to Y_0$ soit de type fini.

